\documentclass[11pt]{article}

\usepackage[margin=1in]{geometry}
\usepackage[T1]{fontenc}
\usepackage{lmodern}

\usepackage{amsmath, amssymb}

\usepackage{setspace}

\usepackage[most]{tcolorbox}
\usepackage{xcolor}

\tcbset{
  enhanced,
  boxrule=0pt,
  frame hidden,
  breakable
}

\usepackage{fancyhdr}
\pagestyle{fancy}
\fancyhf{}
\setlength{\headheight}{52pt}
\renewcommand{\headrulewidth}{0pt}

\newcommand{\Course}{Math 4610 Spring 2026}
\newcommand{\Instructor}{Homework 4}
\newcommand{\AssignmentType}{Homework} % or Discussion
\newcommand{\AssignmentNumber}{4}

\newcommand{\Contributors}{
Marks, Stephen
}

\fancyhead[L]{\Course\\ \Instructor}
\fancyhead[R]{\Contributors}

\newcommand{\ProblemDivider}[1]{%
  \vspace{1.0em}
  \noindent\textbf{#1}\par
  \vspace{0.35em}
}

\definecolor{problemBack}{RGB}{235,242,250}
\definecolor{problemFrame}{RGB}{70,110,160}
\definecolor{exerciseGray}{gray}{0.96}
\definecolor{exerciseGrayFrame}{gray}{0.45}

\newtcolorbox{problemBox}{
  colback=exerciseGray,
  colframe=exerciseGrayFrame,
  arc=2mm,
  left=8pt,right=8pt,top=8pt,bottom=8pt
}

\newtcolorbox{solutionBox}{
  colback=white,
  colframe=white,
  arc=0mm,
  left=0pt,right=0pt,top=0pt,bottom=0pt
}

\newcommand{\ProblemAnnounce}{\noindent}
\newcommand{\SolutionAnnounce}{\noindent\textbf{Solution.}\par\medskip}

\newcommand{\Exercise}[1]{\textbf{Exercise~#1.}\;}

\newcommand{\R}{\mathbb{R}}
\newcommand{\C}{\mathbb{C}}
\newcommand{\F}{\mathbb{F}}
\newcommand{\calL}{\mathcal{L}}

\begin{document}

\begin{center}
  {\Large \textbf{\AssignmentType~\AssignmentNumber}}
\end{center}
\vspace{0.75em}

\ProblemDivider{1}

\begin{problemBox}
\ProblemAnnounce
\Exercise{5.C.6}
Suppose $\F = \C$, $V$ is finite-dimensional, and $T \in \calL(V)$. Prove that if
$k \in \{1, \dots, \dim V\}$, then $V$ has a $k$-dimensional subspace invariant under $T$.
\end{problemBox}

\begin{solutionBox}
\SolutionAnnounce
\doublespacing

Since $V$ is a finite-dimensional complex vector space, every operator $T$ on $V$ has an upper triangular matrix by 5.47. Let $v_1, \dots, v_n$ be a basis of $V$ with respect to which $T$ has an upper triangular matrix. For any $k \in \{1, \dots, n\}$, consider the subspace $U = \operatorname{span}(v_1, \dots, v_k)$. By the conditions for upper triangular matrices (5.39), for each $j \leq k$, we have
\[
Tv_j \in \operatorname{span}(v_1, \dots, v_j) \subseteq \operatorname{span}(v_1, \dots, v_k) = U.
\]
Thus $T(U) \subseteq U$, so $U$ is invariant under $T$. Since $\dim U = k$, we have shown that $V$ has a $k$-dimensional subspace invariant under $T$ for each $k \in \{1, \dots, \dim V\}$.

\end{solutionBox}

\ProblemDivider{2}

\begin{problemBox}
\ProblemAnnounce
\Exercise{5.C.7}
Suppose $V$ is finite-dimensional, $T \in \calL(V)$, and $v \in V$.

(a) Prove that there exists a unique monic polynomial $p_v$ of smallest degree such that $p_v(T)v = 0$.

(b) Prove that the minimal polynomial of $T$ is a polynomial multiple of $p_v$.
\end{problemBox}

\begin{solutionBox}
\SolutionAnnounce
\doublespacing

(a) If $v = 0$, then we can take $p_v = 1$, so assume $v \neq 0$.

Consider the list
\[
v,\, Tv,\, T^2v,\, \dots,\, T^{\dim V}v.
\]
This list contains $1 + \dim V$ vectors in a vector space of dimension $\dim V$, and therefore it must be linearly dependent. By the Linear Dependence Lemma (2.19), there exists a smallest integer
\[
k \in \{1,\dots,\dim V\}
\]
such that
\[
T^k v \in U := \operatorname{span}(v, Tv, \dots, T^{k-1}v),
\]
and such that the list
\[
v, Tv, \dots, T^{k-1}v
\]
is linearly independent. In particular, $\dim U = k$.

We now show that $U$ is invariant under $T$. For $j < k-1$, we have
\[
T(T^j v) = T^{j+1}v \in U,
\]
and by construction,
\[
T(T^{k-1}v) = T^k v \in U.
\]
Thus $T(U) \subseteq U$, so $U$ is invariant under $T$.

Let $p_v$ denote the minimal polynomial of the restricted operator $T|_U$. Because $v \in U$, we have
\[
p_v(T)v = p_v(T|_U)v = 0.
\]

Now suppose $q$ is a polynomial with
\[
\deg q = s < \deg p_v \le \dim U = k
\]
such that
\[
q(T)v = 0.
\]
Then this gives a linear relation expressing $T^s v$ as a linear combination of
\[
v, Tv, \dots, T^{s-1}v,
\]
which contradicts the minimal choice of $k$. Therefore $\deg p_v$ is minimal among all polynomials satisfying $p(T)v = 0$.

Finally, suppose $r$ is another monic polynomial such that
\[
\deg r = \deg p_v
\quad \text{and} \quad
r(T)v = 0.
\]
Then
\[
(p_v - r)(T)v = 0,
\]
and because both $p_v$ and $r$ are monic,
\[
\deg(p_v - r) < \deg p_v.
\]
If $p_v - r$ were not the zero polynomial, we could divide it by its leading coefficient to obtain a monic polynomial $s$ satisfying
\[
s(T)v = 0
\quad \text{and} \quad
\deg s < \deg p_v,
\]
which contradicts the minimality of $\deg p_v$. Thus $p_v - r = 0$, and therefore $p_v = r$.

Thus $p_v$ is unique.

(b) By 5.31, the minimal polynomial of $T$ annihilates every vector in $V$, so $p(T)v = 0$. Since $p_v$ is the minimal degree monic polynomial satisfying $p_v(T)v = 0$, the minimal polynomial of $T$ must be a polynomial multiple of $p_v$.

\end{solutionBox}

\newpage

\ProblemDivider{3}

\begin{problemBox}
\ProblemAnnounce
\Exercise{5.C.8}
Suppose $V$ is finite-dimensional, $T \in \calL(V)$, and there exists a nonzero vector $v \in V$ such that $T^2 v + 2Tv = -2v$.

(a) Prove that if $\F = \R$, then there does not exist a basis of $V$ with respect to which $T$ has an upper-triangular matrix.

(b) Prove that if $\F = \C$ and $A$ is an upper-triangular matrix that equals the matrix of $T$ with respect to some basis of $V$, then $-1 + i$ or $-1 - i$ appears on the diagonal of $A$.
\end{problemBox}

\begin{solutionBox}
\SolutionAnnounce
\doublespacing

(a) The given equation $T^2v + 2Tv = -2v$ can be rewritten as
\[
T^2v + 2Tv + 2v = 0 \quad \Longrightarrow \quad (T^2 + 2T + 2I)v = 0.
\]
Since $v \neq 0$, the operator $T^2 + 2T + 2I$ is not injective. The polynomial $x^2 + 2x + 2$ has roots given by the quadratic formula:
\[
x = \frac{-2 \pm \sqrt{2^2 - 4(1)(2)}}{2(1)} = \frac{-2 \pm \sqrt{4-8}}{2} = \frac{-2 \pm \sqrt{-4}}{2} = \frac{-2 \pm 2i}{2} = -1 \pm i.
\]
Thus $x^2 + 2x + 2 = (x - (-1+i))(x - (-1-i))$, and these roots are non-real complex numbers.

If $\F = \R$, then $T$ has no real eigenvalues (since the only possible eigenvalues would be roots of this polynomial, which are non-real). By 5.44, an operator on a real vector space has an upper-triangular matrix with respect to some basis if and only if its minimal polynomial can be split into linear factors over $\R$. Since the polynomial $x^2 + 2x + 2$ does not split over $\R$, and this polynomial annihilates $v$, the minimal polynomial of $T$ cannot split over $\R$ becuase the minimal polynomial must be a polynomial multiple of the given polynomial by exercise 5.C.7. Therefore, there does not exist a basis of $V$ with respect to which $T$ has an upper-triangular matrix if $\F$ = $\R$.

(b) If $\F = \C$, then by 5.41, the diagonal entries of any upper-triangular matrix representing $T$ are exactly the eigenvalues of $T$. So $(T^2 + 2T + 2I)v = 0$ with $v \neq 0$ implies that $T$ has $-1+i$ or $-1-i$ as an eigenvalue. Thus, either $-1+i$ or $-1-i$ must appear on the diagonal of any upper-triangular matrix representing $T$.

\end{solutionBox}
\newpage

\ProblemDivider{4}

\begin{problemBox}
\ProblemAnnounce
\Exercise{5.C.9}
Suppose $B$ is a square matrix with complex entries. Prove that there exists an invertible square matrix $A$ with complex entries such that $A^{-1}BA$ is an upper-triangular matrix.
\end{problemBox}

\begin{solutionBox}
\SolutionAnnounce
\doublespacing

Suppose $B$ is an $n \times n$ matrix with complex entries. Define a linear operator 
$T \in \calL(\C^n)$ by
\[
T(x) = Bx,
\]
where vectors in $\C^n$ are column vectors.

With respect to the standard basis $e_1, \dots, e_n$ of $\C^n$, the matrix 
of $T$ is exactly $B$.

Because $\F$ = $\C$, 5.47 guarantees the existence 
of a basis $v_1, \dots, v_n$ of $\C^n$ such that the matrix 
$\mathcal{M}(T, (v_1,\dots,v_n))$ is upper triangular.

Let
\[
A = \mathcal{M}(I, (v_1,\dots,v_n), (e_1,\dots,e_n)),
\]
which is the change of basis matrix from $(v_1,\dots,v_n)$ to the standard basis by 3.84. Since 
$(v_1,\dots,v_n)$ is a basis, $A$ is invertible.

By the change of basis formula,
\[
A^{-1}BA
=
\mathcal{M}(T, (v_1,\dots,v_n)),
\]
which is upper triangular. Therefore, there exists an invertible matrix $A$ such 
that $A^{-1}BA$ is upper triangular.

\end{solutionBox}

\end{document}
